% ============================================================================
% PART II: PHYSICS ORGANIZED BY OBSERVABLE FAMILY
% Section 5: Topological Response in 2+1 Dimensions
% ============================================================================
%
% Purpose
%   Fractional Quantum Hall Effect as the canonical laboratory for Chern-Simons
%   TQFT. Transport, fractional charge, braiding, edge modes, and ground-state
%   degeneracy all come from a single observable package (Final QFT Project
%   Overview.md). This section also supplies the Maxwell-Chern-Simons contrast,
%   i.e. how to tell a genuine TQFT from "CS term added to dynamical gauge
%   theory."
%
% Status: PARTIALLY DRAFTED. The Laughlin / quasihole / edge / K-matrix prose
%   below was migrated from archive_2026-05-02_pre_simplification/
%   08_physical_applications.tex. Still TODO: add \cite{...} calls where marked,
%   write the Maxwell-CS subsection (5.5), and write the local-caveats
%   subsection (5.6).
%
% Length target: ~25--30 pages (this is the heart of Part II)
%
% Subsection plan (from Final QFT Project Overview.md)
%   5.1  Hall response
%   5.2  Fractional charge and braiding
%   5.3  Edge modes, anomaly inflow, and ground-state degeneracy
%   5.4  K-matrix theory and nonabelian states
%   5.5  Maxwell--Chern--Simons: when a CS term is not a TQFT
%   5.6  Caveats: mesoscopic complications vs. ideal TQFT observables
%
% Review sources (cited as backbone)
%   Wen 1995 (\cite{Wen1995TopologicalOrdersEdgeExcitations})
%   Stern 2008 (\cite{Stern2008AnyonsQHEReview})
%   Nayak et al. 2008 (\cite{NayakSimonSternFreedmanDasSarma2008NonabelianAnyons})
%   Halperin-Stern-Neder-Rosenow 2011 (\cite{HalperinSternNederRosenow2011FractionalCharge})
%   Halperin-Jain 2020 volume (\cite{HalperinJain2020FQHENewDevelopments})
%   Tong QHE lectures (\cite{TongQuantumHallLectures})
%   Witten 2016 (\cite{Witten2016ThreeLecturesTopologicalPhases})
%
% Landmark theory
%   Laughlin 1983 (\cite{Laughlin1983AnomalousQHE})
%   Arovas-Schrieffer-Wilczek 1984 (\cite{ArovasSchriefferWilczek1984FractionalStatistics})
%   Zhang-Hansson-Kivelson 1989 (\cite{ZhangHanssonKivelson1989FQHE})
%   Wen-Zee 1992 (\cite{WenZee1992AbelianQHClassification})
%   Wen-Niu 1990 (\cite{WenNiu1990GroundStateDegeneracy})
%   Wen 1990 (\cite{Wen1990TopologicalOrders})
%
% Landmark experiments
%   Tsui-Stormer-Gossard 1982 (\cite{TsuiStormerGossard1982FQHEDiscovery})
%   Saminadayar et al. 1997 (\cite{Saminadayar1997FractionalCharge})
%   de-Picciotto et al. 1997 (\cite{dePicciotto1997FractionalCharge})
%   Camino-Zhou-Goldman 2005 (\cite{CaminoZhouGoldman2005ABSuperperiod})
%   Nakamura et al. 2020 (\cite{Nakamura2020AnyonicBraiding})
%   Bartolomei et al. 2020 (\cite{Bartolomei2020AnyonCollisions})
%
% Caveats sources
%   Venkatachalam et al. (\cite{VenkatachalamHartPfeifferWestYacobyGrapheneInterferometers})
%   Rosenow-Halperin (\cite{RosenowHalperinInterferometryPlaceholder})
%
% Writing notes
%   - Subsection 5.6 is the local "iffiness" subsection for the QHE story
%     (from Final QFT Project Overview.md).
%   - Unified four-category taxonomy of strict vs effective observables lives
%     in Sec. 7.2, not here.

\section{Topological Response in $2+1$ Dimensions}
\label{sec:response-2plus1d}

Take a two-dimensional electron gas, place it in a strong perpendicular magnetic field, and cool it to around $T\approx 1\,\text{K}$. Measured across the sample, the Hall conductivity shows the expected integer plateaus at $\nu=1,2,3,\ldots$, but also a staircase of fractional plateaus,
\begin{equation}
  \sigma_{xy}
  =
  \frac{e^2}{h}\,\nu,
  \qquad
  \nu = \frac{1}{3},\,\frac{2}{5},\,\frac{5}{2},\,\ldots
  \label{eq:fqhe-conductance}
\end{equation}
first seen in \cite{TsuiStormerGossard1982FQHEDiscovery}.
The integer case is already understood: at $\nu=n$ one fills $n$ Landau levels, and single-particle quantum mechanics is enough. The fractional plateaus are the harder story. Laughlin's variational wavefunction at $\nu=1/m$ has remarkable numerical overlap with the true ground state \cite{Laughlin1983AnomalousQHE}, but on its own it does not explain why transport is quantized to such extraordinary precision, why the quasihole excitations carry charge $e/m$, or why exchanging two such quasiholes gives a phase $e^{i\pi/m}$ intermediate between bosons and fermions. Each of these is a universal number, and each is observed experimentally. The question is what structure is rigid enough to produce them.

The answer, developed through the late 1980s by Zhang, Hansson, Kivelson, Wen, Niu, Arovas, Schrieffer, Wilczek, and others \cite{ZhangHanssonKivelson1989FQHE,ArovasSchriefferWilczek1984FractionalStatistics,Wen1990TopologicalOrders,WenNiu1990GroundStateDegeneracy}, is that the low-energy physics at $\nu=1/m$ is a Chern--Simons topological field theory. Once that effective theory is written down, the three universal numbers---conductivity, quasihole charge, and anyonic exchange phase---are controlled by a single integer coefficient, and the mathematical structures developed in Part I (differential forms, fiber bundles, linking, functorial quantization, anomaly inflow) are the tools that make the derivation transparent. The rest of this section carries out that derivation in four passes: write the effective action and read off the Hall response, insert quasihole sources and read off the charge and statistics, turn on global topology and read off edge modes and ground-state degeneracy, and finally sketch how the Laughlin case sits inside a broader family. Review treatments that follow the same logic include \cite{Wen1995TopologicalOrdersEdgeExcitations,Stern2008AnyonsQHEReview,NayakSimonSternFreedmanDasSarma2008NonabelianAnyons,HalperinSternNederRosenow2011FractionalCharge,TongQuantumHallLectures,HalperinJain2020FQHENewDevelopments}.

\subsection{Hall Response}
\label{subsec:hall-response}

We cannot solve the microscopic problem---$\sim 10^{11}$ electrons in a wildly degenerate Landau level with Coulomb interactions---and we do not need to. The FQHE state is gapped, so at energies below the gap the correct description is a low-energy effective theory. As usual, we write down the most general action consistent with the symmetries and allowed topology and let experiment fix the coefficients.

For the Laughlin $\nu=1/m$ state, that action \cite{ZhangHanssonKivelson1989FQHE} is
\begin{equation}
  S[a;A]
  =
  \int\left[
  \frac{1}{2\pi}\,A\wedge\dd a
  -
  \frac{m}{4\pi}\,a\wedge\dd a
  \right].
  \label{eq:fqhe-action}
\end{equation}
Here $A$ is the electromagnetic gauge field, treated as a fixed background probe, and $a$ is an \emph{emergent} $U(1)$ gauge field: a low-energy collective mode that packages the correlations of the underlying electrons, in roughly the same sense that a phonon packages the correlations of a crystal lattice. There is no elementary matter field coupled to $a$, and there is no Maxwell term $\int F\wedge\star F$: the bulk theory is purely topological, with no metric and no propagating modes. The mixed term $\frac{1}{2\pi}A\wedge\dd a$ couples the emergent gauge field to electromagnetism, and the pure Chern--Simons term $-\frac{m}{4\pi}a\wedge\dd a$ supplies the dynamics. The integer $m$ is the one parameter the effective theory exposes to us; all three universal numbers follow from it.

The Hall response falls out almost immediately. Treating $a$ as dynamical and varying the action gives the equation of motion
\begin{equation}
  \frac{1}{2\pi}\,\dd A - \frac{m}{2\pi}\,\dd a = 0
  \quad\Longrightarrow\quad
  \dd a = \tfrac{1}{m}\,\dd A,
  \label{eq:fqhe-eom}
\end{equation}
so on any simply connected patch we can take $a = A/m$. Substituting back, the action reduces to
\begin{equation}
  S_{\mathrm{eff}}[A]
  =
  \frac{1}{4\pi m}\int A\wedge\dd A,
  \label{eq:fqhe-effective-action}
\end{equation}
a pure Chern--Simons term for the electromagnetic background. The induced current is
\begin{equation}
  j^\mu
  = \frac{\delta S_{\mathrm{eff}}}{\delta A_\mu}
  = \frac{1}{2\pi m}\,\epsilon^{\mu\nu\rho}\partial_\nu A_\rho,
  \qquad
  j^i = \frac{1}{2\pi m}\,\epsilon^{ij}E_j,
  \label{eq:fqhe-current}
\end{equation}
giving
\begin{equation}
  \sigma_{xy} = \frac{1}{2\pi m} = \frac{e^2}{h}\,\frac{1}{m}
  \label{eq:fqhe-conductance-derived}
\end{equation}
in natural units (restoring $e,\hbar$ in the last step). The fact that the effective action is metric-free is what makes this conductivity a sharp prediction: smooth deformations of the sample geometry, disorder, and local perturbations cannot change a topological invariant continuously. Any change in $\sigma_{xy}$ requires either closing the bulk gap or changing the underlying topological order---i.e.\ a phase transition \cite{Wen1990TopologicalOrders,Wen1995TopologicalOrdersEdgeExcitations}.

We have so far treated the coefficient $m$ in \eqref{eq:fqhe-action} as a continuous parameter. In fact gauge invariance forces it to be an integer, by a large-gauge-transformation argument that is standard by now (see e.g.\ \cite{Dunne1999AspectsCS,TongQuantumHallLectures}). Place the theory on $S^{2}\times S^{1}$, with a unit background monopole flux $\frac{1}{2\pi}\int_{S^{2}}F=1$ through the spatial sphere and Euclidean time compactified on $S^{1}$ of length $\beta$. A large gauge transformation winding once around $S^{1}$ shifts the constant temporal component by $A_0\to A_0+2\pi/\beta$; evaluated on this configuration, the Chern--Simons action shifts by $\Delta S = 2\pi m$. Requiring $\exp(iS)$ to be invariant then forces $m\in\mathbb{Z}$. For electrons, Fermi statistics additionally restrict $m$ to be odd, recovering the observed Laughlin filling fractions $\nu=1,1/3,1/5,\ldots$. The quantization of $\sigma_{xy}$ is therefore built into the effective theory at the level of consistency, not fitted after the fact.

\subsection{Fractional Charge and Braiding}
\label{subsec:fractional-charge-braiding}

The same action determines the excitations. A quasihole worldline is represented by a point source for the emergent gauge field,
\begin{equation}
  S_{\mathrm{source}} = \int a\wedge J,
  \qquad
  J = 2\pi\,\delta^{(2)}(\mathbf{x})\,\dd t
  \label{eq:fqhe-source}
\end{equation}
for a static quasihole at the origin. With this source turned on, the equation of motion \eqref{eq:fqhe-eom} picks up a delta function on its right-hand side, and the emergent magnetic flux through the plane satisfies
\begin{equation}
  m\,f_{12}(\mathbf{x}) = 2\pi\,\delta^{(2)}(\mathbf{x}),
  \label{eq:fqhe-flux-attachment}
\end{equation}
so each unit source drags along $2\pi/m$ units of emergent flux. This is flux attachment. The mixed term in \eqref{eq:fqhe-action} then converts that emergent flux into electric charge density: reading $j^0 = f_{12}/(2\pi)$ from varying with respect to $A_0$, and integrating over space, the quasihole carries
\begin{equation}
  Q_{\mathrm{qh}} = \frac{e}{m}.
  \label{eq:fqhe-fractional-charge}
\end{equation}
The electron has not split into pieces. The Chern--Simons term implements a duality in which the source couples to $a$ via its worldline and then emerges in the electromagnetic current through the mixed term; the ratio, $1/m$, is the same integer that sets the Hall conductivity. One coefficient, two observables. The operational meaning of this fractional charge, and how it has been extracted from shot-noise \cite{Saminadayar1997FractionalCharge,dePicciotto1997FractionalCharge} and interferometry \cite{CaminoZhouGoldman2005ABSuperperiod} measurements, is surveyed in \cite{HalperinSternNederRosenow2011FractionalCharge,Stern2008AnyonsQHEReview}.

The third universal number is the exchange phase of two quasiholes, first derived in \cite{ArovasSchriefferWilczek1984FractionalStatistics}, and it is computed by the same machinery. Two separated quasiholes define two Wilson loops in the emergent gauge field. One quasihole carries unit $a$-charge; the other carries $2\pi/m$ units of emergent flux by \eqref{eq:fqhe-flux-attachment}. The Aharonov--Bohm phase accumulated by one quasihole traversing a closed loop around the other is therefore
\begin{equation}
  \exp\!\Big(\frac{2\pi i}{m}\Big).
  \label{eq:fqhe-full-winding}
\end{equation}
A full winding is two exchanges, so the statistical angle picked up in a single exchange is
\begin{equation}
  \theta = \frac{\pi}{m}.
  \label{eq:fqhe-statistical-angle}
\end{equation}
For $m=1$ this is $\theta=\pi$ (fermions), recovering the integer quantum Hall case; for $m=3$ it is $\theta=\pi/3$. Laughlin quasiholes obey the braid group rather than the permutation group, and the Chern--Simons description ties this exchange phase to the same $m$ that controls $\sigma_{xy}$ and $Q_{\mathrm{qh}}$. Shot-noise, tunnelling, and interferometry measurements all see these fractional quantum numbers in the laboratory \cite{Saminadayar1997FractionalCharge,dePicciotto1997FractionalCharge,CaminoZhouGoldman2005ABSuperperiod,Nakamura2020AnyonicBraiding,Bartolomei2020AnyonCollisions}, which is what makes \eqref{eq:fqhe-action} a physical theory rather than a formal construction.

Both of these computations are special cases of the general Chern--Simons observable theory developed in Sec.~\ref{sec:tqft-observables}. The quasihole charge came from the Gaussian path integral with an external source; the exchange phase came from evaluating a product of Wilson loops, whose expectation value in abelian Chern--Simons is a linking-number phase. The linking-number identity,
\begin{equation}
  \bigl\langle W_{\gamma_1}W_{\gamma_2}\bigr\rangle_{U(1)_m}
  =
  \exp\!\Big(\frac{2\pi i}{m}\,\mathrm{Lk}(\gamma_1,\gamma_2)\Big),
\end{equation}
derived in the earlier observable-theoretic section \cite{Witten1989Jones,Dunne1999AspectsCS}, is exactly what the quasihole braid computes when the two loops are linked once.

\subsection{Edge Modes, Anomaly Inflow, and Ground-State Degeneracy}
\label{subsec:edges-degeneracy}

So far we have worked locally. Two of the most striking features of the FQHE appear only once we ask about the global topology of spacetime: the chiral edge modes on a manifold with boundary, and the topology-dependent ground-state degeneracy on a closed spatial surface \cite{Wen1995TopologicalOrdersEdgeExcitations,WenNiu1990GroundStateDegeneracy}.

The edge story is the Chern--Simons incarnation of anomaly inflow. On a manifold $M$ with boundary $\partial M$, varying the bulk Chern--Simons action produces
\begin{equation}
  \delta S_{\mathrm{CS}}
  =
  \frac{m}{2\pi}\int_{M}F\wedge\delta A
  \;-\;
  \frac{m}{4\pi}\int_{\partial M}A\wedge\delta A.
  \label{eq:fqhe-boundary-variation}
\end{equation}
The bulk term vanishes on solutions of the equations of motion, but the boundary piece survives, and a gauge transformation $A\to A+\dd\lambda$ shifts the action by a boundary integral proportional to $\int_{\partial M}\lambda\,\dd A$. The combined bulk-plus-boundary system can only be gauge invariant if the boundary itself carries a theory whose anomalous variation cancels this inflow. For abelian $U(1)_m$ Chern--Simons the required boundary theory is a chiral boson \cite{Wen1995TopologicalOrdersEdgeExcitations}: solving the bulk flatness condition near $\partial M$ by $A_i = \partial_i\phi$ and substituting back leaves
\begin{equation}
  S_{\partial}
  =
  \frac{m}{4\pi}\int_{\partial M}
  \bigl(\partial_t\phi\,\partial_x\phi
  - v\,(\partial_x\phi)^2\bigr)\,\dd t\,\dd x,
  \qquad
  (\partial_t + v\partial_x)\phi = 0,
  \label{eq:fqhe-chiral-boson}
\end{equation}
with $v$ a nonuniversal edge velocity. The boundary field propagates in only one direction, so the edge theory is chiral; its universal features, including the chiral central charge, are fixed by the bulk topological data. For the Laughlin $\nu=1/m$ states this is a single chiral boson with $c=1$, so the quantized thermal Hall conductance is $\kappa_{xy} = c\,\pi^{2}k_{B}^{2}T/(3h) = \pi^{2}k_{B}^{2}T/(3h)$, independent of $m$. Thermal transport, shot-noise, and tunnelling experiments on fractional edges all see behaviour of this form, and the detailed spectrum of the edge agrees with what \eqref{eq:fqhe-chiral-boson} predicts \cite{Wen1995TopologicalOrdersEdgeExcitations,Witten2016ThreeLecturesTopologicalPhases}. From this viewpoint the edge mode is not an additional phenomenological boundary ingredient; it is the degree of freedom required to restore gauge invariance of the combined bulk-boundary system, which in turn is what anomaly inflow already told us to expect.

The second global consequence is the ground-state degeneracy \cite{WenNiu1990GroundStateDegeneracy}. For a closed spatial surface $\Sigma$, the equations of motion set $F=0$, so we are quantizing the moduli space of flat $U(1)$ connections as in Sec.~\ref{subsec:hilbert-surfaces-gsd}. On a torus this space is parameterized by the two holonomies $u=\oint_\alpha A$, $v=\oint_\beta A$ with $u,v\sim u+2\pi,v+2\pi$, and the Chern--Simons action reduces to a symplectic quantum-mechanics problem on this $T^{2}_{\mathrm{hol}}$ with symplectic form
\begin{equation}
  \Omega = \frac{m}{2\pi}\,\dd u\wedge\dd v,
  \qquad
  \int_{T^{2}_{\mathrm{hol}}}\Omega = 2\pi m.
  \label{eq:fqhe-symplectic}
\end{equation}
Geometric quantization assigns one quantum state per $2\pi\hbar$ unit of phase-space volume \cite{ElitzurMooreSchwimmerSeiberg1989CanonicalCS}, giving
\begin{equation}
  \dim\mathcal{H}_{U(1)_m}(T^{2}) = m.
  \label{eq:fqhe-torus-degeneracy}
\end{equation}
For $\nu=1/3$, the torus ground-state manifold is three-dimensional, and threading external flux through one of the cycles rotates the system through those three sectors with period $h/e$; that period is the experimental signature that the degeneracy is topological rather than an accidental consequence of a symmetry. The torus result extends naturally to higher-genus surfaces, where the Hilbert-space dimension becomes $m^{g}$ on genus $g$ \cite{WenNiu1990GroundStateDegeneracy}. This is the abelian $U(1)_m$ specialization of the Verlinde formula, which the functorial-quantization discussion in Sec.~\ref{subsec:hilbert-surfaces-gsd} assembled from a pair-of-pants decomposition and which reappears in WZW conformal blocks, theta functions on abelian varieties, and integrable representations of affine Lie algebras at level $m$. The same integer that fixes $\sigma_{xy}$ and $\theta$ in the bulk fixes how many ground states the topology admits.

\subsection{$K$-Matrix Theory and Nonabelian States}
\label{subsec:k-matrix-nonabelian}

Most observed filling fractions are not Laughlin. States like $\nu=2/5$ and $\nu=5/2$ require more than a single emergent gauge field, and the natural generalization is the $K$-matrix framework \cite{WenZee1992AbelianQHClassification}. Introduce $N$ emergent $U(1)$ gauge fields $a^I$ and a charge vector $\mathbf{t}\in\mathbb{Z}^N$, and take the action
\begin{equation}
  S_{K}[a^I;A]
  =
  \int\left[
  \frac{1}{4\pi}K_{IJ}\,a^I\wedge\dd a^J
  +
  \frac{1}{2\pi}t_I\,A\wedge\dd a^I
  \right],
  \label{eq:fqhe-k-matrix-action}
\end{equation}
with $K$ an integer symmetric matrix; the Laughlin action \eqref{eq:fqhe-action} is the $1\times 1$ case $K=(m)$, $\mathbf{t}=(1)$. Integrating out the $a^I$ gives
\begin{equation}
  S_{\mathrm{eff}}[A]
  =
  \frac{\mathbf{t}^{\mathsf{T}}K^{-1}\mathbf{t}}{4\pi}\int A\wedge\dd A
  \quad\Longrightarrow\quad
  \sigma_{xy} = \frac{e^2}{h}\,\mathbf{t}^{\mathsf{T}}K^{-1}\mathbf{t},
  \label{eq:fqhe-k-hall}
\end{equation}
while quasiparticles, labelled by integer vectors $\boldsymbol{\ell}\in\mathbb{Z}^N$, carry charges and mutual statistics
\begin{equation}
  Q_{\boldsymbol{\ell}} = e\,\mathbf{t}^{\mathsf{T}}K^{-1}\boldsymbol{\ell},
  \qquad
  \theta_{\boldsymbol{\ell},\boldsymbol{\ell}'} = 2\pi\,\boldsymbol{\ell}^{\mathsf{T}}K^{-1}\boldsymbol{\ell}',
\end{equation}
and the genus-$g$ ground-state degeneracy is $|\det K|^{g}$. The Halperin $(3,3,1)$ state is a minimal illustration:
\begin{equation}
  K = \begin{pmatrix} 3 & 1 \\ 1 & 3\end{pmatrix},
  \qquad
  \mathbf{t} = \begin{pmatrix}1\\1\end{pmatrix},
  \qquad
  K^{-1} = \frac{1}{8}\begin{pmatrix} 3 & -1 \\ -1 & 3\end{pmatrix},
\end{equation}
from which $\sigma_{xy} = (e^{2}/h)\cdot 2/5$, the elementary quasiparticles $(1,0)$ and $(0,1)$ each carry charge $e/4$, and their mutual statistical angle is $2\pi\cdot 1/8 = \pi/4$. Every abelian FQHE state fits into this framework \cite{WenZee1992AbelianQHClassification,Wen1995TopologicalOrdersEdgeExcitations,HalperinJain2020FQHENewDevelopments}.

The abelian story does not exhaust the phenomenology. The $\nu=5/2$ plateau, observed in ultra-clean samples in the second Landau level, is the most striking example that does not: it is widely believed to be described by the Moore--Read (Pfaffian) state, which is a \emph{nonabelian} Chern--Simons theory, specifically $SU(2)_{2}$ \cite{NayakSimonSternFreedmanDasSarma2008NonabelianAnyons}. The elementary quasiparticle is a spinor $\sigma$ with fusion rule $\sigma\times\sigma = 1 + \psi$, so a collection of $n$ such quasiparticles spans a Hilbert space of dimension $\sim 2^{(n-2)/2}$; exchanging them executes unitary transformations on this space rather than multiplying it by a phase. We will not pursue the nonabelian construction here. For present purposes it is enough to note that the Laughlin states are not isolated examples but the simplest, abelian cases of a much larger structure: for every compact Lie group $G$ and level $k$, Chern--Simons theory $G_k$ assigns Hilbert spaces to surfaces, Wilson-loop expectation values to knots, and boundary conformal field theories to edges \cite{Witten1989Jones,Dunne1999AspectsCS}, and the FQHE realizes several of these theories in the laboratory. If the $\nu=5/2$ identification is confirmed, nonabelian braiding in such a system would furnish the first fault-tolerant topological qubit in a solid-state platform \cite{NayakSimonSternFreedmanDasSarma2008NonabelianAnyons}.

\subsection{Maxwell--Chern--Simons: When a CS Term Is Not a TQFT}
\label{subsec:maxwell-cs}

% TODO: Write this subsection.
% Source material: archive_2026-05-01/anomalies_boundaries_topological_response.tex
% lines ~500-600 (to be recovered from the archive via a migration pass).
%
% Content outline:
%   - Action S = -(1/4 e^2) int F wedge *F + (k/4pi) int A wedge dA.
%   - The Hodge star *F reintroduces metric; propagator has a topologically
%     massive pole m_CS = k e^2 / (2 pi). See \cite{Dunne1999AspectsCS}.
%   - Even though the theory has a CS term, it is NOT a TQFT: observables
%     depend on the metric, and there are local propagating modes.
%   - Useful as an effective theory: finite-thickness edges, finite-
%     temperature regularization of the pure-CS limit.
%   - Lesson: the purity of the CS action matters. Adding a CS term to an
%     otherwise dynamical gauge theory gives a topologically massive gauge
%     theory, not a topological one.
% Cite: \cite{Dunne1999AspectsCS}, \cite{Wen1995TopologicalOrdersEdgeExcitations},
%       \cite{BirminghamBlauRakowskiThompson1991TopologicalFieldTheory}

\subsection{When the Observables Are Clean, and When They Are Only Effective}
\label{subsec:qhe-caveats}

% TODO: Write this subsection.
% Role: local "iffiness" discussion for the 2+1d story
% (Final QFT Project Overview.md). The broader four-part taxonomy across all
% observable families (strict TQFT / IR response / theory-side / model
% dependent) lives in Sec.~7.2.
%
% Bullets (from tqft_observables_literature_review.md Sec. 4.5):
%   - Interferometry in real devices suffers Coulomb-dominated regimes, area
%     fluctuations, slow quasiparticle dynamics, edge reconstruction; the
%     clean TQFT phase is not the mesoscopic device observable.
%   - Edge-mode counting: chirality and anomaly coefficient are topological;
%     velocities, equilibration lengths, tunneling exponents are not.
%   - The K-matrix does not encode all geometric / gravitational response;
%     the Wen-Zee term and gravitational CS terms are additional data.
% Cite: \cite{VenkatachalamHartPfeifferWestYacobyGrapheneInterferometers},
%       \cite{RosenowHalperinInterferometryPlaceholder},
%       \cite{Wen1995TopologicalOrdersEdgeExcitations},
%       \cite{Stern2008AnyonsQHEReview}

\medskip

Looked at from far enough away, the FQHE story is simple. Once the low-energy effective theory is fixed by the symmetries of the problem to be a Chern--Simons theory, a single integer controls the quantized transport, the fractionalization of quasiparticle charge, the braiding phase, the chiral edge, and the topology-dependent ground-state degeneracy---and gauge invariance forces that integer to exist. The mathematics developed earlier---differential forms, bundles, topological actions, functorial quantization, and anomaly inflow---is the language in which that rigidity becomes transparent. In that sense the fractional quantum Hall effect is not only an application of topological field theory, but the experimental evidence that topological field theory governs real macroscopic matter.
